\documentclass[11pt,a4paper]{article}

\usepackage[T1]{fontenc}
\usepackage[utf8]{inputenc}
\usepackage{lmodern}
\usepackage{textcomp}

\usepackage[
  top=2.5cm, bottom=3.0cm,
  left=3.0cm, right=2.8cm,
  headheight=15pt
]{geometry}

\usepackage{setspace}
\setlength{\parindent}{0pt}
\setlength{\parskip}{0.55em}
\linespread{1.12}

% ---- Colors ----
\usepackage{xcolor}
\definecolor{primary}{HTML}{FF5F5F}
\definecolor{secondary}{HTML}{2D4A7A}
\definecolor{accent}{HTML}{C44D4D}
\definecolor{dark}{RGB}{30, 30, 35}
\definecolor{light}{RGB}{252, 248, 248}
\definecolor{codebg}{HTML}{F5F0F0}
\definecolor{mebox}{HTML}{FFF0F0}

% ---- Packages ----
\usepackage{tcolorbox}
\tcbuselibrary{breakable}
\usepackage{enumitem}
\usepackage{listings}
\usepackage{booktabs}
\usepackage{tabularx}

\lstset{
  basicstyle=\small\ttfamily\color{dark},
  backgroundcolor=\color{codebg},
  frame=none,
  breaklines=true,
  xleftmargin=1em,
  framexleftmargin=1em,
}

% ---- ME box ----
\newtcolorbox{mebox}{
  colback=mebox,
  colframe=primary!70,
  coltitle=dark,
  arc=4pt,
  boxrule=0.8pt,
  left=8pt, right=8pt, top=6pt, bottom=6pt,
  breakable,
  before upper={\bfseries},
}

% ---- Headers ----
\usepackage{fancyhdr}
\pagestyle{fancy}
\fancyhf{}
\fancyhead[L]{\small\color{secondary} Ideation Chat Log}
\fancyfoot[R]{\small\thepage}
\renewcommand{\headrulewidth}{0.4pt}
\renewcommand{\footrulewidth}{0pt}

\usepackage[
  colorlinks=true,
  linkcolor=secondary,
  urlcolor=secondary,
  citecolor=secondary,
]{hyperref}

\title{\color{primary}\textbf{Ideation Chat Log}\\[0.3em]
\large\color{secondary} Welfare Discovery \& Rule Engine Research}
\author{}
\date{}

\begin{document}
\maketitle
\thispagestyle{fancy}

\bigskip
\hrule
\bigskip

% ---- Turn 1 ----

{\small\color{primary}\textbf{dpg}}
\begin{mebox}
India has lot of welfare schemes but many are ambiguous and hard to discover. Are there companies that simplify this for everyone?
\end{mebox}

{\small\color{secondary}\textbf{Gemini}}\vspace{2pt}

Navigating the maze of Indian welfare schemes is famously difficult, but a growing ecosystem of ``Civic Tech'' and ``Social Impact'' companies is dedicated to solving exactly this. By 2026, these players have moved beyond simple lists to using AI-driven eligibility engines and ``last-mile'' agents to ensure benefits actually reach people.

Here are the key companies and platforms simplifying discovery today:

\subsection*{1. The Market Leaders}
\begin{itemize}[nosep]
\item \textbf{Haqdarshak:} The most prominent company in this space. They operate a tech platform that maps thousands of central and state schemes. Their unique model uses a network of trained ``Haqdarshaks'' (women entrepreneurs in villages) who use an app to screen citizens for eligibility and help with application paperwork for a small fee.
\item \textbf{myScheme (Govt of India):} A government initiative under Digital India functioning like a modern startup platform. A ``one-stop search engine'' where you enter your age, caste, state, and income, and it filters to show only schemes you qualify for.
\end{itemize}

\subsection*{2. Specialized \& Sector-Specific Players}
\begin{itemize}[nosep]
\item \textbf{Indus Action:} A high-impact non-profit focusing on specific ``vulnerable moments'' (birth, school admission, job loss). Famous for their work on RTE section 12(1)(c), helping thousands of low-income families navigate private school admissions.
\item \textbf{Rocket Learning:} AI-powered WhatsApp bots helping parents in low-income communities discover early childhood education and nutrition schemes.
\item \textbf{Civis India:} Focuses on the ``pre-legislation'' stage, simplifying complex draft laws and welfare policies into plain language.
\end{itemize}

\subsection*{3. AI \& Language-First Solutions}
\begin{itemize}[nosep]
\item \textbf{Kisan e-Mitra:} AI chatbot for farmers to discover agricultural subsidies using voice commands in their mother tongue.
\item \textbf{Suman Sakhi:} Healthcare-focused AI helping rural women discover maternal and newborn welfare benefits through WhatsApp.
\end{itemize}

\subsection*{4. Niche \& Business Focused}
\begin{itemize}[nosep]
\item \textbf{Startup India Portal:} Aggregates 200,000+ recognized startups and matches them with tax exemptions, seed funds, and patent subsidies.
\item \textbf{Benani:} Newer fintechs integrating ``scheme discovery'' into their apps, checking if transaction history suggests eligibility for schemes like PM Shram Yogi Maan-dhan.
\end{itemize}

\textbf{How to use them:}
\begin{enumerate}[nosep]
\item \textbf{Direct Search:} Go to myScheme.gov.in for a quick eligibility check.
\item \textbf{Assisted Help:} Download the Haqdarshak app for help with documentation and last-mile submission.
\end{enumerate}

\bigskip
\hrule
\bigskip

% ---- Turn 2 ----

{\small\color{primary}\textbf{dpg}}
\begin{mebox}
Any of these provide API access for free? and have they solved the problem per say. What is it that lacks in this space of welfare-discovery
\end{mebox}

{\small\color{secondary}\textbf{Gemini}}\vspace{2pt}

Whether any of these provide free API access depends on how you define ``free'' and what layer of the stack you want to access. As of 2026, the ecosystem is split between government infrastructure, open-source protocols, and commercial platforms.

\subsection*{1. Can you get free API access?}
The short answer is: \textbf{Yes, but it's tiered.}

\begin{itemize}[nosep]
\item \textbf{myScheme \& API Setu (The Government Route):} Connected via API Setu. Government departments and registered ``Authorized Service Providers'' can access these APIs for free or at nominal cost.
\item \textbf{Indus Action (The Open-Source Route):} They have pioneered Scheme Markup Language (SML)---an open-source protocol turning complex government text into machine-readable rules.
\item \textbf{Haqdarshak (The Commercial/B2B Route):} Generally no free public API. They sell their API/SaaS stack to state governments, banks, and large NGOs.
\end{itemize}

\subsection*{2. Have they ``solved'' the problem?}
Technically? \textbf{Almost.} Logistically? \textbf{No.}

We have solved the \textbf{Information Gap}: Between myScheme and Haqdarshak, nearly 5,000+ schemes are now digitized and searchable.

However, they haven't solved the \textbf{Fulfillment Gap}. Knowing you are eligible is only 10\% of the battle:
\begin{itemize}[nosep]
\item \textbf{The ``Broken Link'' Problem:} You find a scheme on a discovery app, but the link leads to a state portal that is down, or a PDF from 2021.
\item \textbf{The Document Loop:} You have to upload your Aadhaar, Income Certificate, and Caste Certificate 10 separate times to 10 different departments.
\end{itemize}

\subsection*{3. What is still lacking?}
\begin{itemize}[nosep]
\item \textbf{Interoperability (The ``Apply Once'' Dream):} No ``Universal Shopping Cart'' for welfare.
\item \textbf{Proactive Eligibility (Push vs.\ Pull):} Currently the citizen must search. The holy grail is automatic notification and pre-approval.
\item \textbf{Dynamic Rule Matching:} Schemes change criteria frequently. Many tools struggle to keep rule-engines updated in real-time across 28 states.
\item \textbf{The ``Grievance'' Black Hole:} Discovery apps are powerless when an application is stuck in ``pending'' for six months.
\end{itemize}

\textbf{In short:} We've built the ``Google'' of welfare (Discovery); now we need the ``Amazon'' of welfare (one-click checkout and tracking).

\bigskip
\hrule
\bigskip

% ---- Turn 3 ----

{\small\color{primary}\textbf{dpg}}
\begin{mebox}
Wait but how is just adding a layer of tracking/applying once and connect digilocker for the apply once feature not implemented in the existing commercial solutions? Are you sure these are clearly a space untapped?

Also what do you mean by Dynamic Rule matching? What I understand is you mean we have a proprietary rule engine but that needs to be updated everytime smthg changes, instead one can have a middleware that does this dynamically without HITL?
\end{mebox}

{\small\color{secondary}\textbf{Gemini}}\vspace{2pt}

If it were just a technical problem of ``connecting APIs,'' it would have been solved years ago. The reason we don't have ``Single-Click Apply'' for all 5,000+ Indian welfare schemes is less about technology and more about \textbf{sovereignty, liability, and the ``API-gap.''}

\subsection*{1. Why ``Apply Once'' isn't fully solved}
\begin{itemize}[nosep]
\item \textbf{Document vs.\ Workflow APIs:} DigiLocker provides Document APIs (a PDF of your Aadhaar). To ``apply once,'' you need a Workflow API for the specific department. Most government portals are ``Walled Gardens.''
\item \textbf{The ``Authorized Agent'' Liability:} Legally, applying for a government scheme is a ``declaration of truth.'' The DPDP Act 2023 makes it very difficult for a private company to ``sign'' on behalf of a citizen without fresh, explicit OTP/Consent for every application.
\item \textbf{The State-Silo Problem:} 28 states, each with its own tech stack. Many state schemes aren't ``API-fied'' for submission.
\item \textbf{The ``Liveness'' Requirement:} Many schemes require Aadhaar-based Face Auth or a thumbprint at a physical kiosk (CSC).
\end{itemize}

\subsection*{2. Dynamic Rule Matching: The Middleware Dream}
Historically, companies had a Hard-Coded Rule Engine. If the government raised the income limit, a human programmer had to change \texttt{if (income <= 200000)}.

\textbf{What is happening now (The Middleware Shift):} LLM-driven Ingestion using protocols like Indus Action's SML.
\begin{itemize}[nosep]
\item An AI agent monitors the Gazette of India and state government notification PDFs. When a new notification is released, the AI ``reads'' the legal prose and suggests a code update to the rule engine.
\item \textbf{Why not 100\% ``No-HITL''?} In welfare, a 0.1\% Hallucination Rate is a catastrophe---10,000 families wrongly rejected for food rations.
\item \textbf{The Current Standard:} ``AI-First, Human-Verified.'' The AI writes the logic, a legal expert clicks ``Approve.''
\end{itemize}

\subsection*{3. Is the space truly untapped?}
It is \textbf{partially tapped}. The Haryana Model (Parivar Pehchan Patra / Family ID) is the closest to ``Apply Once.'' The state proactively starts your pension when you turn 60.

The opportunity: a B2B Middleware connecting private companies (Zomato, Swiggy) to government APIs for ``Welfare-as-a-Benefit'' for gig workers automatically.

\bigskip
\hrule
\bigskip

% ---- Turn 4 ----

{\small\color{primary}\textbf{dpg}}
\begin{mebox}
So the solution that is actually needed is the country to have a standardised scheme system for each state? Also wdym zomato swiggy to govt apis? What benefits? basic income etc?
\end{mebox}

{\small\color{secondary}\textbf{Gemini}}\vspace{2pt}

The short answer is \textbf{yes}. The ultimate solution isn't one giant ``super-app,'' but a Standardized Protocol---a common language. Think of it like \textbf{UPI for Welfare}. Just as UPI allowed any bank to talk to any payment app, a ``Welfare Protocol'' (currently being built on the Beckn Protocol) allows any government scheme to talk to any platform.

\subsection*{1. What does ``Zomato/Swiggy to Govt APIs'' mean?}
Under the Code on Social Security (2020), gig workers are being brought into the formal safety net. Platforms like Zomato, Swiggy, and Blinkit now connect their backend to e-Shram via APIs. When a delivery partner completes 90 days of work, the platform's API automatically pings the government portal.

\subsection*{2. Is it ``Basic Income''?}
Generally \textbf{not} UBI. Instead, it focuses on Risk Protection:
\begin{itemize}[nosep]
\item Accident \& Disability Insurance
\item Health Cover (Ayushman Bharat PM-JAY, INR 5 Lakh free hospital coverage)
\item Old Age Protection (PM-SYM pension, co-paid by government and platform)
\item Maternity Benefits
\end{itemize}

\subsection*{3. Is ``Standardized Schemes'' the actual solution?}

\begin{tabularx}{\textwidth}{@{}lXX@{}}
\toprule
\textbf{Feature} & \textbf{Old Way (Silos)} & \textbf{New Way (Standardized Protocol)} \\
\midrule
Discovery & Search 50 state portals & One app searches everything \\
Data & Upload Aadhaar 50 times & One-time consent via Account Aggregator \\
Eligibility & Clerk manually checks papers & Rule Engine matches in real-time \\
Application & Wait weeks for ``Update'' & Track via unified API Dashboard \\
\bottomrule
\end{tabularx}

\bigskip
\hrule
\bigskip

% ---- Turn 5 ----

{\small\color{primary}\textbf{dpg}}
\begin{mebox}
Okay so this beckn protocol will be replacing all existing fragmented welfare schemes? isnt this what ondc uses? but how is this commerce protocol helpful here?

Also what other segments of workers/people are not captured in this welfare scheme systems that exist rn. Like gig workers is one of them and?

So a solution that I make should combine good rule engines/easy to navigate/beckn protocol/new customer segments right? Any R\&D on the technical side of things that may help or everything is solved? Because I assume rule engines exist and only refinement at best can be done or am I wrong?
\end{mebox}

{\small\color{secondary}\textbf{Gemini}}\vspace{2pt}

\subsection*{1. Beckn: It's not a ``Commerce'' Protocol, it's a ``Discovery'' Protocol}
ONDC uses Beckn, but Beckn itself is ``sector-agnostic.'' HTTP is the protocol for websites---it doesn't care if the website sells shoes or shows news. In welfare, Beckn helps a citizen ``discover'' a scheme and ``transact'' (apply for) the benefit. It allows ``Unbundling''---a private startup builds the UI, while the government provides the Registry and Fulfillment backend.

\subsection*{2. Who is being left behind?}
\begin{itemize}[nosep]
\item \textbf{Migrant Workers (The Portability Gap):} A worker from Bihar in a factory in Karnataka can't access local state schemes because documents are tied to home state.
\item \textbf{Home-Based Workers:} Millions of women working from home (rolling agarbattis, stitching garments). ``Invisible'' to e-Shram---no formal employer.
\item \textbf{Domestic Workers:} Most maid/cook registries are local and informal.
\item \textbf{The ``Missing Middle'':} People just above the Poverty Line who don't get free rations but are one medical emergency away from poverty.
\end{itemize}

\subsection*{3. The R\&D Frontier: What isn't ``solved'' yet?}

\textbf{A. Policy-as-Code (Converting Legalese to Logic)}\\
Using LLMs to read a 50-page State Gazette, extract eligibility logic, and output standardized JSON schema that any app can use instantly.

\textbf{B. Zero-Knowledge Proofs (ZKP) for Eligibility}\\
A middleware that queries the Income Tax database and only sends a True/False signal to the welfare app. The app knows you are eligible but never sees your actual income data.

\textbf{C. Federated Eligibility Checking}\\
Instead of one central Master Database, federated queries ping Aadhaar API, Farmer Registry API, and Ration Card API simultaneously to ``triangulate'' eligibility without storing data.

\bigskip
\hrule
\bigskip

% ---- Turn 6 ----

{\small\color{primary}\textbf{dpg}}
\begin{mebox}
I like the suggestions you gave but wdym ZKP for eligibility???? How tf can a private startup bring that to life. It's obviously in the hands of the govt to extend that feat.

For A recommendation, I like the json idea but wouldn't that be very cumbersome to debug and audit? What are other simpler ways per se.
\end{mebox}

{\small\color{secondary}\textbf{Gemini}}\vspace{2pt}

\subsection*{1. How a Startup brings ZKP to Life}
A startup cannot ``invent'' your income data. However, in the India Stack's Verifiable Credentials (VC) ecosystem:
\begin{itemize}[nosep]
\item The startup acts as a ``Verifier,'' not an ``Issuer.'' The Government (via API Setu / DigiLocker) is the ``Issuer.''
\item When a user wants to prove they earn $<$INR 2L, they use an Account Aggregator (AA). The AA fetches data, and the verifier logic runs a local check on the user's device.
\item The app generates a cryptographic ``Proof'' saying ``This user's income is $<$INR 2L, signed by the Govt's public key.'' Only this proof is passed to the welfare department's API.
\item The startup's value: building the ``Trust Interface'' that manages these ZKP handshakes.
\end{itemize}

\subsection*{2. Simpler alternatives to JSON for Welfare Rules}

\textbf{DMN (Decision Model and Notation):}\\
Decision Tables---like an Excel sheet on steroids. A human can read the rows: \texttt{If Age > 60 AND Income < 3L -> Eligibility = TRUE}. Converts to execution logic via Drools or DMN.js.

\textbf{DSL (Domain Specific Language):}\\
A human-readable language for welfare:
\begin{lstlisting}
REQUIRE age >= 18
REQUIRE state == 'Maharashtra'
GRANT benefit 'Rs. 5000'
\end{lstlisting}
Much easier to debug than JSON because it reads like English but executes like code.

\textbf{LLM-Verified Auditing:}\\
Feed your rule engine to an LLM and ask: ``Summarize all edge cases where a widow in rural Karnataka would be rejected.'' The LLM audits logic and identifies bugs in plain language.

\bigskip
\hrule
\bigskip

% ---- Turn 7 ----

{\small\color{primary}\textbf{dpg}}
\begin{mebox}
The indus action SML u told abt is not present anywhere
\end{mebox}

{\small\color{secondary}\textbf{Gemini}}\vspace{2pt}

Apologies---Scheme Markup Language (SML) isn't a ``website'' or a single downloadable app. It's a protocol---a set of rules for how to write code so different systems can talk to each other about welfare.

Indus Action treats SML as part of their ``Digital Public Good'' stack. You won't find a standalone repo called ``SML,'' but it is the logic backbone of their Eligibility Engines. It is officially listed under their ``Protocols that Scale'' on the Indus Action Model page. In the wild, look into their GitHub Organization (Indus-Action-Initiatives) at the RTE-MIS or Eligibility Engine repositories.

\textbf{The gap for you is to build the ``IDE for Welfare'':} A platform where a non-technical person can type a welfare policy in plain English, and your middleware converts it into a Beckn-compliant API with ZKP privacy baked in.

\end{document}
